\scanpage{87}
\mspage{74}

Nous allons présenter les choses de façon un peu plus générale.

\uline{Définition} 4.4. J'appelle structure de \unsure{cofibrations|cofibration} \uline{spéciale} une structure $(M,W,C)$ satisfaisant les conditions a) à f) de 4.1., \deleted{\unsure{que}} je \unsure{reprends|reformule}

\textcircled{a} $M$ stable par petites limites inductives, et les objets de $M$ sont accessibles,

\textcircled{b} $W$ est \unsure{fbt|ftt|fort} saturée, (\deleted{\unsure{de}} $\mathrm{Cof}_0$ de p. 39), $C$ stable par cochgt. de base et par compositions \deleted{\uline{transfinies},} $\mathrm{TC}\overset{\mathrm{déf}}{=}C\cap W$ stable par cochgt. de base (et par composition automatiquement)

\textcircled{b$'$} $C$ et $\mathrm{TC}$ sont stables $\unsure{=}$ par composition \uline{transfinie}

\textcircled{c} $\exists$ \uline{petite} partie $C_0\subset C$ telle que $C\subset\widetilde C_0$ ($\Longleftrightarrow\widetilde C=\widetilde C_0$)

\textcircled{d} $\unsure{\deleted{C}|{}}\mathrm{TF}\overset{\mathrm{déf}}{=}C_*\subset W$ (ce qui implique, \inserted{joint à \textcircled{a}--\textcircled{c},} l'axiome de factorisation $\mathrm{Cof}_3$ de page 39).

4.5. Nous allons donner une recette générale d'extension d'une telle structure, quand on dispose d'un couple de foncteurs adjoints
\diagram{\begin{tikzcd}[ampersand replacement=\&,column sep=huge] M \arrow[r,shift left,"\unsure{u|v}_!"] \& M' \arrow[l,shift left,"u^*"]
\end{tikzcd}\,,}
$M$ étant muni d'une structure de cofibration spéciale \inserted{$(W,C)$}, et $M'$ satisfaisant la condition liminaire \textcircled{a}. On \unsure{définira|définit} alors les parties $W',\deleted{\unsure{(C)}}C'$ de $\operatorname{Fl}M'$
% archive sheet number: 86
